\subsection{一致空间的完备化}

定理 3 设 X 是一个一致空间。则存在完备 Hausdorff 一致空间 \(\hat{\mathbf{X}}\) 与一致连续映射 \(i: \mathbf{X} \to \hat{\mathbf{X}}\)，它们具有下列性质：

\begin{enumerate}
\def\labelenumi{(\Alph{enumi})}
\setcounter{enumi}{15}
\tightlist
\item
  给定任意一致连续映射 \(f\)，它把 \(\mathbf{X}\) 映入完备 Hausdorff 一致空间 \(\mathbf{Y}\)，则存在唯一的一致连续映射 \(g: \hat{\mathbf{X}} \to \mathbf{Y}\) 使 \(f = g \circ i\)。
\end{enumerate}

若 \((i_1, X_1)\) 是由完备 Hausdorff 一致空间 \(X_1\) 与具有性质 (P) 的一致连续映射 \(i_1: X \to X_1\) 组成的另一个偶，则存在唯一的同构 \(\varphi: \hat{X} \to X_1\) 使 \(i_1 = \varphi \circ i\)。

定理的第一个断言表明，偶 \((i, \hat{\mathbf{X}})\) 是下列泛映射问题（《集合论》，第四章，§ 3，第 1 目）的解：其中的 \(\Sigma\)-集是完备 Hausdorff 一致空间，\(\sigma\)-态射是一致连续映射，而 \(\alpha\)-映射是 X 到完备 Hausdorff 一致空间中的一致连续映射。于是，偶 \((i, \hat{\mathbf{X}})\) 在相差一个唯一同构的意义下的唯一性由泛映射问题之解的一般性质（同处所引）得出。剩下要证明的是偶 \((i, \hat{\mathbf{X}})\) 的存在性。

\begin{enumerate}
\def\labelenumi{\arabic{enumi})}
\tightlist
\item
  \(\hat{\mathbf{X}}\) 的定义。令 \(\hat{\mathbf{X}}\) 为 \(\mathbf{X}\) 上（第 2 目）极小 Cauchy 滤子全体组成的集。我们要在 \(\hat{\mathbf{X}}\) 上定义一个一致结构。为此，设 \(\mathbf{V}\) 是 \(\mathbf{X}\) 的任一对称一致邻域，令 \(\tilde{\mathbf{V}}\) 表示具有一个公共 V-小集的极小 Cauchy 滤子偶 \((\mathfrak{X}, \mathfrak{Y})\) 全体组成的集。我们将证明，诸集合 \(\tilde{\mathbf{V}}\) 构成 \(\hat{\mathbf{X}}\) 上某个一致结构的一致邻域基本系：
\end{enumerate}

\begin{enumerate}
\def\labelenumi{(\roman{enumi})}
\item
  由于每个 \(\mathcal{X} \in \hat{\mathbf{X}}\) 都是 Cauchy 滤子，按定义我们有 \((\mathcal{X}, \mathcal{X}) \in \tilde{\mathbf{V}}\)，其中 \(\mathbf{V}\) 取遍 \(\mathbf{X}\) 的所有对称一致邻域；因此公理 \((\mathbf{U}_{\mathbf{I}}^{\prime})\) 成立。
\item
  若 V 与 \(V'\) 是 X 的两个对称一致邻域，则 \(W = V \cap V'\) 是对称一致邻域，且每个 W-小的集合也是 V-小的与 \(V'\)-小的；因此 \(\tilde{W} \subset \tilde{V} \cap \tilde{V}'\)，这就证明了 \((B_{I})\)。
\item
  诸集合 \(\tilde{\mathbf{V}}\) 按定义是对称的，因此 \((\mathbf{U}_{\mathbf{II}}^{\prime})\) 成立。（iv）设 \(\mathbf{V}\) 为 \(\mathbf{X}\) 的一个对称一致邻域，\(\mathbf{W}\) 为使 \(\tilde{\mathbf{V}} \subset \mathbf{V}\) 的一个对称一致邻域。考虑三个极小 Cauchy 滤子 \(\mathcal{X}, \mathcal{Y}, \mathfrak{Z}\)，使得 \((\mathcal{X}, \mathcal{Y}) \in \tilde{\mathbf{W}}\) 且 \((\mathcal{Y}, \mathfrak{Z}) \in \tilde{\mathbf{W}}\)；于是存在两个 W-小集 \(\mathbf{M}, \mathbf{N}\)，使 \(\mathbf{M} \in \mathcal{X} \cap \mathcal{Y}\) 且 \(\mathbf{N} \in \mathcal{Y} \cap \mathfrak{Z}\)。由于 \(\mathbf{M}\) 与 \(\mathbf{N}\) 都属于 \(\mathcal{Y}\)，\(\mathbf{M} \cap \mathbf{N}\) 非空，因此（第 1 目命题 1）\(\mathbf{M} \cup \mathbf{N}\) 是 \(\tilde{\mathbf{W}}\)-小的，从而是 V-小的；由于 \(\mathbf{M} \cup \mathbf{N}\) 既属于 \(\mathcal{X}\) 又属于 \(\mathfrak{Z}\)，我们有 \(\tilde{\mathbf{W}} \subset \tilde{\mathbf{V}}\)；因此 \((\mathbf{U}_{\mathbf{III}}^{\prime})\) 成立。
\end{enumerate}

我们接着证明一致空间 \(\hat{\mathbf{X}}\) 是 Hausdorff 的。设 \(\mathcal{X},\mathcal{Y}\) 是 \(\mathbf{X}\) 上两个极小 Cauchy 滤子，使得 \((\mathcal{X},\mathcal{Y})\in \hat{\mathbf{X}}\) 对 \(\mathbf{V}\) 取遍 \(\mathbf{X}\) 的所有对称一致邻域都成立。由此立即得出，诸集合 \(\mathbf{M}\cup \mathbf{N}\)（其中 \(\mathbf{M}\in \mathcal{X}\)、\(\mathbf{N}\in \mathcal{Y}\)）构成一个滤子 \(\mathfrak{Z}\) 的基，该滤子比 \(\mathcal{X}\) 与 \(\mathcal{Y}\) 都粗。现在 \(\mathfrak{Z}\) 是 Cauchy 滤子，因为当 \(\mathbf{V}\) 取遍 \(\mathbf{X}\) 的对称一致邻域时，由假设总存在一个 V-小集 \(\mathbf{P}\)，它同时属于 \(\mathcal{X}\) 与 \(\mathcal{Y}\)，因而属于 \(\mathfrak{Z}\)。由极小 Cauchy 滤子的定义，我们有 \(\mathcal{X} = \mathfrak{Z} = \mathcal{Y}\)，这就表明 \(\hat{\mathbf{X}}\) 是 Hausdorff 的。

\begin{enumerate}
\def\labelenumi{\arabic{enumi})}
\setcounter{enumi}{1}
\item
  \(i\) 的定义；\(\mathbf{X}\) 的一致结构是经由 \(i\) 从 \(\hat{\mathbf{X}}\) 的一致结构取逆像而得的。我们知道，对每个 \(x \in \mathbf{X}\)，邻域滤子 \(\mathfrak{B}(x)\)（\(x\) 在 \(\mathbf{X}\) 中的邻域滤子）是极小 Cauchy 滤子（第 2 目命题 5 推论 1）。于是我们定义 \(i(x) = \mathfrak{B}(x)\)。令 \(f = i \times i\)；我们将证明，对 \(\mathbf{V}\) 取遍 \(\mathbf{X}\) 的每个对称一致邻域，都有 \(\overline{j}(\tilde{\mathbf{V}}) \subset \mathbf{V} \cup \overline{j}[(\mathbf{V})^{\sim}]\)，而这将证明我们的断言（§ 2 第 4 目）。现在，若 \([i(x), i(y)] \in \tilde{\mathbf{V}}\)，则存在一个 V-小集 \(\mathbf{M}\)，它既是 \(x\) 的邻域又是 \(y\) 的邻域，因此 \((x, y) \in \mathbf{V}\)。反之，若 \((x, y) \in \mathbf{V}\)，则立即可见集合 \(\mathbf{V}(x) \cup \mathbf{V}(y)\) 是 \(\overline{\mathbf{V}}\)-小的，且既是 \(x\) 的邻域又是 \(y\) 的邻域。
\item
  \(\hat{\mathbf{X}}\) 是完备的且 \(i(\mathbf{X})\) 在 \(\hat{\mathbf{X}}\) 中稠密。在 \(i(\mathbf{X})\) 上，邻域 \(\tilde{\mathbf{V}}(\mathfrak{X})\)（它是点 \(\mathfrak{X} \in \hat{\mathbf{X}}\) 的邻域）的迹，是使下列关系成立的那些 \(i(x)\) 组成的集：
\end{enumerate}

\[
(\mathcal {X}, i (x)) \in \tilde {\mathbf {V}}.
\]

这个关系意味着，存在 \(x\) 在 X 中的一个 V-小邻域，它属于 \(\mathcal{X}\)，也就是说，\(x\) 是 \(\mathcal{X}\) 的某个 V-小集的内点。令 M 为 \(\mathcal{X}\) 的所有 V-小集的内部之并；则 M 属于 \(\mathcal{X}\)（第 2 目命题 5 推论 4），并且由刚才所说的可得 \(\tilde{\mathrm{V}}(\mathcal{X}) \cap i(\mathrm{X}) = i(\mathrm{M})\)。我们由此断言：

\begin{enumerate}
\def\labelenumi{(\roman{enumi})}
\item
  \(\tilde{\mathbf{V}}(\mathcal{X}) \cap i(\mathbf{X})\) 非空，因此 \(i(\mathbf{X})\) 在 \(\hat{\mathbf{X}}\) 中稠密。
\item
  \(\tilde{\mathbf{V}}(\mathcal{X})\) 在 \(i(\mathbf{X})\) 上的迹属于 X 上的滤子基 \(i(\mathcal{X})\)；因此这个滤子基在 \(\hat{X}\) 中收敛于点 \(\mathfrak{X}\)。
\end{enumerate}

现在设 \(\mathfrak{F}\) 是 \(i(\mathbf{X})\) 上的一个 Cauchy 滤子；则由上面的 2) 及第 1 目命题 4，\(\overline{i}^{1}(\mathfrak{F})\) 是某个 Cauchy 滤子 \(\mathfrak{G}\)（\(\mathbf{X}\) 上的）的基。设 \(\mathfrak{X}\) 是比 \(\mathfrak{G}\) 粗的极小 Cauchy 滤子（第 2 目命题 5）；则 \(i(\mathfrak{X})\) 是 \(i(\mathbf{X})\) 上的一个 Cauchy 滤子基（第 1 目命题 3），且 \(\mathfrak{F}=i[\overline{i}^{1}(\mathfrak{F})]\) 比以 \(i(\mathfrak{X})\) 为基的滤子细。由于后者在 \(\hat{\mathbf{X}}\) 中收敛，\(\mathfrak{F}\) 亦然，于是第 4 目命题 9 表明 \(\hat{\mathbf{X}}\) 是完备的。

\begin{enumerate}
\def\labelenumi{\arabic{enumi})}
\setcounter{enumi}{3}
\tightlist
\item
  性质 (P) 的验证。设 \(f\) 是 \(\mathbf{X}\) 到完备 Hausdorff 一致空间 \(\mathbf{Y}\) 中的一致连续映射。我们先证明，存在唯一的一致连续映射 \(g_{0}: i(\mathbf{X}) \to \mathbf{Y}\) 使 \(f = g_{0} \circ i\)。由于 \(f\) 连续，我们有
\end{enumerate}

\[
f (x) = \lim f (\mathfrak {B} (x)),
\]

因此，若令 \(g_0(i(x)) = \lim f(\mathfrak{V}(x))\)，则我们有 \(f = g_0 \circ i\)；于是剩下要证明 \(g_0\) 在 \(i(\mathbf{X})\) 上一致连续。设 \(\mathbf{U}\) 是 \(\mathbf{Y}\) 的一个一致邻域，\(\mathbf{V}\) 是 \(\mathbf{X}\) 的一个对称一致邻域，使得关系 \((x, x') \in \mathbf{V}\) 蕴含 \((f(x), f(x')) \in \mathbf{U}\)；我们在 2) 中已看到，关系 \((i(x), i(x')) \in \tilde{\mathbf{V}}\) 蕴含 \((x, x') \in \mathbf{V}\)，因而也蕴含

\[
(g _ {0} (i (x)), g _ {0} (i (x ^ {\prime}))) \in \mathbf {U},
\]

这就证明了我们的断言。

设 \(g\) 是 \(g_0\) 由连续性到 \(\hat{\mathbf{X}}\) 上的延拓（第 6 目定理 2）；则 \(f = g \circ i\)，并且显然 \(g\) 是 \(\hat{\mathbf{X}}\) 到 \(Y\) 中满足此关系的唯一连续映射，因为 \(i(\mathbf{X})\) 在 \(\hat{\mathbf{X}}\) 中稠密（第一章 § 8 第 1 目命题 2 推论 1）。

证毕。

定义 4 在定理 3 的证明中所定义的完备 Hausdorff 一致空间 \(\hat{\mathbf{X}}\) 称为 \(\mathbf{X}\) 的 Hausdorff 完备化，而映射 \(i: \mathbf{X} \to \hat{\mathbf{X}}\) 称为 \(\mathbf{X}\) 到其 Hausdorff 完备化中的典范映射。

我们还指出下列事实：

命题 12 （i）子空间 \(i(\mathbf{X})\) 在 \(\hat{\mathbf{X}}\) 中稠密。

\begin{enumerate}
\def\labelenumi{(\roman{enumi})}
\setcounter{enumi}{1}
\item
  等价关系 \(i(x) = i(x')\) 的图像是 \(\mathbf{X}\) 的诸一致邻域之交。
\item
  X 的一致结构是 \(\hat{X}\) 的一致结构在 i 下的逆像{[}或者说是子空间 \(i(\mathbf{X})\) 的一致结构在 i 下的逆像{]}。
\item
  \(i(\mathbf{X})\) 的一致邻域是 X 的一致邻域在 \(i \times i\) 下的像，而在 \(\hat{X} \times \hat{X}\) 中，\(i(\mathbf{X})\) 的一致邻域的闭包构成 \(\hat{X}\) 的一致邻域基本系。
\item
  与 (iii) 已在定理 3 的证明过程中证明；（iv）则由 (i) 与 (iii) 借助于先前证明的一般结果（§ 2 第 4 目的注与命题 6）得出。关系式
\end{enumerate}

\[
i (x) = i \left(x ^ {\prime}\right)
\]

按定义意味着 \(x\) 与 \(x'\) 有相同的邻域滤子。但按定义这就蕴含 \((x, x') \in V\)（其中 \(V\) 是 \(X\) 的任意一致邻域），而其逆是显然的。

推论 若 X 是 Hausdorff 一致空间，则典范映射 \(i: X \to \hat{X}\) 是 X 到 \(\hat{X}\) 的一个稠密子空间上的同构。

当 X 是 Hausdorff 空间时，称 \(\hat{X}\) 为 X 的完备化，并且我们通常借助 i 把 X 与 \(\hat{X}\) 的一个稠密子集等同起来。

注。若作了这一等同，则 X 上的极小 Cauchy 滤子恰好是 \(\hat{X}\) 的点的邻域滤子在 X 上的迹；这由定理 3 的证明得出。

命题 12 的推论刻画了 Hausdorff 一致空间的完备化：

命题 13 若 Y 是完备 Hausdorff 一致空间，且 X 是 Y 的稠密子空间，则典范单射 \(\mathbf{X} \to \mathbf{Y}\) 可延拓为 \(\hat{\mathbf{X}}\) 到 Y 上的一个同构。

因为由第 6 目定理 2，X 到完备 Hausdorff 一致空间 Z 中的每个一致连续映射都唯一地延拓为 Y 到 Z 中的一致连续映射。

命题 14 设 X 是完备 Hausdorff 一致空间，U 是它的一致结构，并设 Z 是 X 的稠密子空间。若 U' 是 X 上的一个比 U 粗、且在 Z 上诱导出与 U 相同的一致结构的一致结构，则 U = U'。

设 \(X'\) 表示集合 \(X\) 赋以一致结构 \(\mathcal{U}'\) 后所得的一致空间。典范映射 \(X' \to \hat{X}'\) 与恒同映射 \(X \to X'\) 的复合是一致连续映射 \(\varphi: X \to \hat{X}'\)。由于 \(Z\) 关于由 \(\mathcal{U}'\) 诱导的一致结构是 Hausdorff 的，\(\varphi\) 在 \(Z\) 上的限制依假设是 \(Z\) 到 \(\varphi(Z)\)（\(\hat{X}'\) 的稠密子空间）上的同构；于是（第 6 目定理 2 的推论）\(\varphi\) 本身是 \(X\) 到 \(\hat{X}'\) 上的同构，因此 \(X' = \hat{X}'\) 且 \(\mathcal{U}' = \mathcal{U}\)。

命题 15 设 X 与 X' 是两个一致空间。对每个一致连续映射 \(f: \mathbf{X} \to \mathbf{X}'\)，存在唯一的一致连续映射 \(\hat{f}: \hat{\mathbf{X}} \to \hat{\mathbf{X}}'\) 使图表

\[
\begin{array}{c}\mathbf {X} \stackrel {{i ^ {\prime}}} {{\to}} \mathbf {X} ^ {\prime}\\i \downarrow \quad \quad \quad \quad \quad \quad \quad \quad \quad \quad \quad \quad \quad \quad \quad \quad \quad \quad \quad \quad \hat {\mathbf {X}} \stackrel {{\rightarrow}} {{\longrightarrow}} \hat {\mathbf {X}} ^ {\prime}\\f\end{array}
\]

可交换 (*)，其中 \(i: \mathbf{X} \to \hat{\mathbf{X}}\) 与 \(i': \mathbf{X}' \to \hat{\mathbf{X}}'\) 是典范映射。

对函数 \(i' \circ f: \mathbf{X} \to \hat{\mathbf{X}}'\) 应用定理 3。

推论 若 \(f: \mathbf{X} \to \mathbf{X}'\) 与 \(g: \mathbf{X}' \to \mathbf{X}''\) 是两个一致连续映射，且 \(h = g \circ f\)，则 \(\hat{h} = \hat{g} \circ \hat{f}\)。

这是命题 15 中唯一性的直接推论。

